\subsection{Data Requirements}

PWC requires tract-level population data, not just district polygon geometry.
This distinguishes PWC from PP, Reock, CH, S, and LW, which depend only on
the geometric shape of the district polygon.
PWC computation requires the 2020 Census PL 94-171 redistricting file,
which provides:
\begin{itemize}
  \item Total population $p_i$ for each tract $i$.
  \item Tract geographic centroid $c_i$ (internal point or area centroid from TIGER/Line).
\end{itemize}

\subsection{Centroid Computation}

Tract centroids are the geographic centroids of tract polygons,
computed as the area-weighted mean of polygon vertices.
We use the ``interior point'' method: for non-convex tracts, we use the
point-on-surface (POS) rather than the area centroid to ensure the centroid
is inside the tract.
The area centroid of a non-convex tract can fall outside the tract itself,
producing misleading distance computations for tracts with holes or inlets.

\subsection{Analytical Example}

Consider a $3 \times 3$ grid of unit tracts with equal population $p = 1$.
Tract centroids are at $(i + 0.5, j + 0.5)$ for $i, j \in \{0, 1, 2\}$.
The population-weighted centroid of the whole grid is:
\[
  \bar{c} = \frac{\sum_{i,j} (i + 0.5, j + 0.5)}{9} = (1.5, 1.5).
\]
PWC for the whole-grid district:
\begin{align}
  \text{PWC} &= \frac{1}{9} \sum_{i,j} \|(i+0.5, j+0.5) - (1.5, 1.5)\|^2 \\
  &= \frac{1}{9} \sum_{i,j} \left[(i-1)^2 + (j-1)^2\right] \\
  &= \frac{1}{9} \cdot 9 \cdot \frac{2(0^2 + 1^2 + 1^2)}{3}
  = \frac{4}{3} \approx 1.333 \text{ square units}.
\end{align}
This matches the L1 test invariant: PWC equals the sum of squared distances
from each tract centroid to the grid centre, divided by 9.

\subsection{Implementation}

Implemented in \texttt{crates/bisect-analysis/src/compactness.rs}.
Requires the run's tract assignment file and the Census tract centroid file.
PWC for each district is computed in $O(m)$ where $m$ is the number of tracts
in the district.
Total computation for all districts is $O(M)$ where $M$ is the total tract count.
